\section{Introduction}
\label{sec:introduction}

Compression asks which distinctions in a source should survive a constrained representation and which can be discarded.
In a reconstruction problem, this choice is expressed through the cost of reproducing the source inaccurately: a finite description is useful when it preserves the features that matter at an acceptable cost~\citep{gray1998quantization,cover2006elements}.
For many sources, however, the structure of interest lies in relationships among elements.
A graph specifies connections among vertices, a distance or similarity matrix records geometry among observations, and a teacher representation induces affinities among examples.
Preserving these relationships gives a natural compression question: how can their structure be represented more economically while retaining the distinctions selected by a fidelity criterion?

Several seemingly unrelated methods approach this question, some explicitly and others through representation learning.
Graph summaries retain reduced descriptions of adjacency structure~\citep{riondato2017graph}; similarity-preserving hashing represents observations by compact codes~\citep{weiss2008spectral}; and relational distillation transfers relationships from teacher to student representations~\citep{park2019relational,tung2019similarity}.
Their descriptions, fidelity criteria, and constrained resources differ, but they can be examined through the same set of choices: the source relation, the retained description, the relation recovered from it, the fidelity of that recovery, and its resource cost.
Making these choices explicit separates representation, decoder, fidelity, and resource: the same retained assignments can support different relational decoders, while similar nominal capacity or fidelity under one criterion need not imply the same code organization or preservation under another.

We formulate \emph{relational compression} in the classical source--description--reconstruction sense, with relational structure as the fidelity-bearing source content.
For a fixed or shared carrier set, its basic interface is
\begin{equation}
S\longrightarrow m_S\longrightarrow\widetilde S,
\label{eq:relational-interface-intro}
\end{equation}
where $S$ is the source relation, $m_S$ is its retained description, and $\widetilde S$ is the relation reconstructed or evaluated from that description.
A relational distortion specifies how preservation of source structure is assessed, while a resource functional measures the cost of retaining it.
Together, these components turn relational preservation into a constrained compression problem: what relational fidelity can be retained at a given description resource?
Specific instances of this pattern occur across established literatures, including explicit relational and graph-compression formulations discussed in \cref{sec:methods-map}.
The contribution here is to make the source relation, retained description, reconstructed or evaluated relation, fidelity criterion, and resource convention explicit as a common cross-domain interface while retaining the assumptions specific to each construction.

Within this framework, we develop finite-codeword equality as a representation mechanism.
A finite assignment groups elements according to their codewords; probabilistic assignments replace hard equality by the probability of a same-codeword collision.
This gives a setting in which to study how a source relation shapes the representation through alignment and separation, and how aggregate codeword use describes its organization.
The main contributions are:
\begin{itemize}
    \item \textbf{A framework for relational compression.}
    We formalize relational compression as a source--description--reconstruction problem in which relational structure is the fidelity-bearing content.
    The framework explicitly separates the source relation, retained description, reconstructed or evaluated relation, relational fidelity criterion, constrained resource, and shared context.
    This provides a common interface for formulating and comparing relation-preserving methods across domains while retaining the assumptions specific to each construction.

    \item \textbf{Finite-codeword collision as a reusable realization.}
    Within this framework, we develop same-codeword collision as a general representation-side mechanism for finite representations.
    The resulting construction connects pair-specific alignment and separation, aggregate R\'enyi-2 occupancy, reference-weighted affinity, and positive-spherical geometry, and exposes exact correspondences with squared-Euclidean centroid reconstruction and normalized graph association and cut.
    Controlled synthetic, graph, and image studies illustrate how the same collision machinery can be instantiated under different relational and reconstruction-defined fidelity requirements.
\end{itemize}

The remainder of the paper is organized as follows.
\Cref{sec:relational-compression-framework} defines the relational-compression problem and its general fidelity constructions.
\Cref{sec:methods-map} presents the method mappings and worked summary example.
\Cref{sec:collision-geometry,sec:collision-entropy,sec:alignment-separation} develop finite-codeword collision geometry and its objectives, followed by selected realizations in \cref{sec:relational-distortion-realizations} and experimental studies in \cref{sec:experiments}.


\section{A Framework for Relational Compression}
\label{sec:relational-compression-framework}

Relational compression requires identifying the relational structure to preserve, the description retained about it, how that description is decoded or evaluated, and how fidelity is traded against resource.
This section formalizes those components and introduces pairwise, content-based, and operator-based fidelities.
We begin with a simple block-valued summary example to show how they fit together.

Let $V$ be a finite carrier, let $S$ be a pairwise relation on $V$, let $c:V\to\{1,\ldots,K\}$ assign each element to a group, and let $\Xi=(\Xi_{ab})_{a,b=1}^K$ be a $K\times K$ table of reproduction values.
The retained description $m_S=(c,\Xi)$ defines the block-valued relational summary
\begin{equation}
\widetilde S(v,w)=\Xi_{c(v),c(w)}.
\label{eq:block-valued-relational-summary}
\end{equation}
The assignments determine which pairs share a reproduction value, while the table determines the value reproduced for those pairs.
A fidelity criterion measures how well $\widetilde S$ preserves $S$, and a resource measure accounts for the retained assignments and table values.
The same finite assignments could support a different decoder, so they do not by themselves uniquely determine the reconstructed relation.
This is a compression formulation because the full relation is represented through the constrained description $m_S=(c,\Xi)$, trading the resource required to retain that description against the fidelity of the relation reconstructed from it by a decoder.

\subsection{Relational sources}

A relational source consists of elements together with structure defined among them.
For a graph, the carrier elements are vertices and the relation records adjacency or edge weights; for a collection of observations, the carrier elements are the observations and the relation may record pairwise distances, similarities, or neighborhood structure.
Separating the elements from their relation identifies the content whose preservation will be assessed.

Let $V$ be a finite carrier set.
A pairwise relational object is
\begin{equation}
\mathfrak R=(V,S),
\qquad
S:V\times V\to\mathcal Y,
\end{equation}
where $S(v,w)$ is the relation value assigned to the pair $(v,w)$.
The codomain $\mathcal Y$ specifies the type of relation value, and the admissible class $\mathcal S_V\subseteq\mathcal Y^{V\times V}$ specifies its structural properties, such as symmetry or nonnegativity.
Reconstructed relations belong to a class $\widehat{\mathcal S}_V\subseteq\widehat{\mathcal Y}^{V\times V}$.
Thus a reconstructed object is $\widetilde{\mathfrak R}=(V,\widetilde S)$, with the same carrier and a reproduced relation.

We treat the carrier identities as fixed or shared between encoder and decoder.
The compression target is therefore $S$ conditional on $V$.
Additional shared context $C$ may contain a reproduction alphabet, a carrier ordering, or model parameters used across source objects.
The retained description contains the source-specific information required for reconstruction beyond this context.
We suppress $C$ in the notation when it is fixed.

In learned realizations, an element $v$ may be presented to the encoder through an observation $x_v$, with $X_V=(x_v)_{v\in V}$ denoting the collection of inputs.
These observations provide the information from which the retained description is constructed, while the relational structure $S$ specifies the relationships among carrier elements that are to be represented.
For example, a graph encoder may receive node features $x_v$ while the source relation $S$ is the graph adjacency or edge-weight structure.
The same per-object formulation applies to datasets of relational objects $(V_i,S_i)$ with varying carrier sets.

\subsection{Descriptions and reconstruction}

The block-valued summary above already separates the stored description from the relation reconstructed from it.
We now formalize that distinction.

A compression scheme specifies a description space $\mathcal M_V$, an encoder $\mathsf{Enc}_V$, and a decoder $\mathsf{Dec}_V:\mathcal M_V\to\widehat{\mathcal S}_V$.
For a source relation $S$, write
\begin{equation}
m_S=\mathsf{Enc}_V(S),
\qquad
\widetilde S=\mathsf{Dec}_V(m_S).
\label{eq:relational-reconstruction-interface}
\end{equation}
Additional encoder-side observations can be included through $\mathsf{Enc}_V(S,X_V)$.
The description may contain stored relation values, element assignments, coordinates, or parameters of a relational reconstruction rule.
The decoder can materialize $\widetilde S$ or evaluate its entries on demand through $(m_S,v,w)\mapsto\widetilde S(v,w)$.
Either form gives a relation determined by the description and shared context.

With a common source and reconstruction codomain, a scheme is \emph{lossless} on $\mathcal S_V$ when
\begin{equation}
\mathsf{Dec}_V(\mathsf{Enc}_V(S))=S
\qquad
\text{for every }S\in\mathcal S_V.
\label{eq:lossless-relational-compression}
\end{equation}
In the lossy setting, the reconstructed relation approximates the source according to a specified fidelity criterion.
For the block-valued summary in \cref{eq:block-valued-relational-summary}, the table reproduces exactly those source relations that are constant on its evaluated blocks; variation within a block becomes approximation error.
The description and decoder therefore determine the family of relational structures available for reconstruction.

This is the familiar encoder--description--decoder interface of classical compression applied to a relational source~\citep{cover2006elements}.
\Cref{tab:framework-components} summarizes its components together with the fidelity and resource measures introduced next.

\begin{table}[t]
\centering
\small
\setlength{\tabcolsep}{4pt}
\renewcommand{\arraystretch}{1.12}
\setlength{\extrarowheight}{1pt}
\begin{tabularx}{\linewidth}{>{\raggedright\arraybackslash}p{0.18\linewidth}
                            >{\raggedright\arraybackslash}X
                            >{\raggedright\arraybackslash}X}
\toprule
Role & Classical compression & Relational compression \\
\midrule
Source
& Source object $x$
& Relational object $(V,S)$, with $V$ fixed or shared \\
Description
& Retained message $m$
& Retained relational description $m_S$ \\
Reconstruction
& Reproduced object $\widehat x$
& Reconstructed relation $\widetilde S$ \\
Decoder
& Rule mapping $m$ to $\widehat x$
& Rule constructing $\widetilde S$ or evaluating its entries \\
Fidelity
& Distortion $d(x,\widehat x)$
& Relational distortion $\Delta_V(S,\widetilde S)$ \\
Resource
& Cost of the description
& Cost $\Omega(m_S)$ in specified units \\
Shared context
& Information available to both sides
& Carrier set and shared context $C$ \\
\bottomrule
\end{tabularx}
\caption{The compression interface for ordinary and relational source objects.}
\label{tab:framework-components}
\end{table}

Structured descriptions can determine many pair relations from comparatively few retained components; subsequent sections instantiate this principle with block summaries and finite codes.

\subsection{Resources and relational fidelity}
\label{sec:relational-distortion}

Given a description and decoder, compression becomes a trade-off between description cost and relational distortion: the resource measure records the former and the fidelity criterion the latter.

A resource functional
\begin{equation}
\Omega:\mathcal M_V\to\R_{\ge0}
\end{equation}
measures description cost in specified units, such as encoded bits, stored relation entries, or scalar coordinates.
A relational distortion
\begin{equation}
\Delta_V:\mathcal S_V\times\widehat{\mathcal S}_V\to\R_{\ge0}
\end{equation}
measures distortion of the source relation according to the chosen fidelity criterion.
Whenever exact source reconstruction is admissible, we take $\Delta_V(S,S)=0$.
More generally, zero distortion means that the chosen fidelity criterion assigns no distortion to the reconstruction; it need not imply literal equality between $S$ and $\widetilde S$.
The decoder determines $\widetilde S$ from the retained description and any shared context $C$, while fidelity is assessed against the source $S$.
Thus $\Delta_V$ specifies how well the reconstructed or evaluated relation preserves the source according to the declared criterion.
A training objective may equal this distortion, include it as one term, or optimize a surrogate for it.

For a source relation $S$, the best distortion attainable within the declared description space and decoder under a resource budget $B$ is
\begin{equation}
D_V^\star(B;S)
=
\inf_{\substack{m\in\mathcal M_V\\\Omega(m)\le B}}
\Delta_V\bigl(S,\mathsf{Dec}_V(m)\bigr).
\label{eq:lossy-relational-compression}
\end{equation}
An encoder optimized for this criterion selects, or approximates, a minimizing admissible description for each source relation.
As the budget increases, the feasible set expands, so the best attainable distortion is nonincreasing.
Thus $D_V^\star(B;S)$ describes the resource--distortion trade-off for the source: how much relational distortion remains at each available description budget.

The definition of $\Delta_V$ is deliberately general.
We next give several concrete fidelity constructions used in subsequent sections, beginning with a pairwise criterion.

\subsection{Pairwise relational fidelity}

For pairwise relational sources, a natural construction is to evaluate relation values locally and average their distortion over a declared distribution of carrier pairs:
\begin{equation}
\Delta_V(S,\widetilde S)
=
\E_{(\mathsf V,\mathsf V')\sim P_{\mathrm{pair}}}
\bigl[d_{\mathrm{rel}}(S(\mathsf V,\mathsf V'),\widetilde S(\mathsf V,\mathsf V'))\bigr].
\label{eq:pairwise-relational-fidelity}
\end{equation}
Here $P_{\mathrm{pair}}$ is a probability distribution on carrier pairs, $\mathsf V,\mathsf V'$ denote a random pair drawn according to this distribution, and $d_{\mathrm{rel}}:\mathcal Y\times\widehat{\mathcal Y}\to\R_{\ge0}$ is a local relation-value distortion function.
Uniform pair weights with squared error give average relation-value reconstruction error.
Pair distributions concentrated on graph edges, neighboring observations, or selected queries emphasize those portions of the source instead.
Thus $P_{\mathrm{pair}}$ determines which source relations receive weight, while $d_{\mathrm{rel}}$ determines how relational reconstruction error is measured locally.

This construction is useful when relational fidelity is naturally assessed through individual relation values.
Other applications evaluate larger-scale structural properties of the relation; we consider two such constructions next.

\subsection{Selected structural fidelity constructions}

Depending on which properties of the source relation are important, a criterion may instead evaluate retained structural content, the action of the relation on a family of probes, or another source-dependent property.
We develop two related constructions here because they are used in the graph realizations later in the paper, while other source-dependent fidelity constructions remain possible.

\paragraph{Content-based fidelity}

Some fidelity criteria are naturally expressed in terms of relational content retained during compression.
For example, deleting graph edges can remove direct connections and the paths supported by them.
A content functional assigns a value to the structure that remains, allowing fidelity to be measured by the decrease in that value.

Let $\preceq_{\mathrm{rel}}$ be a compatible, composable simplification relation on the family under consideration.
A nonnegative structural-content functional $\mathcal I$ is monotone when
\begin{equation}
S_2\preceq_{\mathrm{rel}}S_1
\quad\Longrightarrow\quad
\mathcal I(S_2)\le\mathcal I(S_1).
\label{eq:relational-content-monotone}
\end{equation}
For $\widetilde S\preceq_{\mathrm{rel}}S$, define
\begin{equation}
D_{\mathcal I}(S,\widetilde S)
=\mathcal I(S)-\mathcal I(\widetilde S).
\label{eq:relational-content-distortion}
\end{equation}
The resulting distortion is nonnegative and increases or stays constant along a sequence of admissible simplifications.
For $\mathcal I(S)>0$, the normalized version is
\begin{equation}
\overline D_{\mathcal I}(S,\widetilde S)
=1-\frac{\mathcal I(\widetilde S)}{\mathcal I(S)}.
\label{eq:normalized-relational-content-distortion}
\end{equation}
It measures the fraction of source content lost: exact retention gives zero, and a null reconstruction with zero content gives one.

Choosing $\mathcal I$ specifies the structural property to be preserved.
For example, total retained edge weight gives direct relation distortion, while path- or walk-based content values connections supported over several steps, and nonnegative weighted sums combine several such properties.
The graph realizations later in the paper use this construction to connect simple relational deletions with their effects on source structure.

\paragraph{Operator-based fidelity}

Given a fixed carrier $V$, we use a relation-derived operator to mean a linear map $T_S:\R^V\to\R^V$ determined by the source relation $S$.
For a graph, familiar examples include adjacency, transition, and Laplacian operators; the symmetric fidelity construction below and the graph realizations later in the paper primarily use the Laplacian.
A relation can also be evaluated through the behavior it supports on its carrier.
For a graph, a signal assigns a value to each vertex, and the graph Laplacian measures how that signal varies across connections~\citep{shuman2013emerging}.
Such signals act as probes: comparing their energy before and after compression measures which aspects of the source geometry have been retained.
This leads to a fidelity construction that will supply several of the graph objectives used later.

Let $T_S$ and $T_{\widetilde S}$ be finite-dimensional symmetric operators on signals over $V$, and suppose
\begin{equation}
0\preceq T_{\widetilde S}\preceq T_S,
\label{eq:operator-compression-order}
\end{equation}
where $\preceq$ is the Loewner order.
For a fixed probe $f$, the removed energy is
\begin{equation}
D_f(S,\widetilde S)
=f^\top(T_S-T_{\widetilde S})f.
\label{eq:operator-probe-distortion}
\end{equation}
The operator order ensures that this quantity is nonnegative.
It measures fidelity along the particular signal direction selected by $f$.

To evaluate a family of directions, let $F$ be a random probe with finite second moment $M=\E[FF^\top]$.
The average distortion is
\begin{equation}
\begin{aligned}
D_M(S,\widetilde S)
&=\E\bigl[F^\top(T_S-T_{\widetilde S})F\bigr]\\
&=\operatorname{tr}\bigl[(T_S-T_{\widetilde S})M\bigr].
\end{aligned}
\label{eq:operator-expected-distortion}
\end{equation}
The matrix $M$ determines how signal directions are weighted.
A rank-one choice $M=ff^\top$ recovers the fixed-probe distortion, while $M=I$ weights coordinate directions equally.
Source-derived choices can emphasize modes of the original relation; these are computed from that source and held fixed while candidate reconstructions are compared.
For fixed $M\succeq0$, the average operator distortion is itself a content-based distortion: defining $\mathcal I_M(S)=\operatorname{tr}(T_SM)$ makes \cref{eq:operator-expected-distortion} exactly the corresponding content difference.
The operator viewpoint remains useful because it provides a structured way to specify which signal directions should be preserved and supports, for example, a worst-case criterion that asks how much relative energy can be lost in any source-supported direction:
\begin{equation}
D_{\mathrm{wc}}(S,\widetilde S)
=\sup_{f:f^\top T_Sf>0}
\frac{f^\top(T_S-T_{\widetilde S})f}{f^\top T_Sf}.
\label{eq:operator-worst-case-distortion}
\end{equation}
Under \cref{eq:operator-compression-order}, this lies in $[0,1]$ for $T_S\ne0$; we set it to zero when $T_S=0$.
The average criterion expresses fidelity under a chosen probe distribution, while the worst-case criterion measures the largest relative distortion.
Spectral sparsification gives a related all-probe graph objective, using reweighted sparse reconstructions to preserve Laplacian quadratic forms across signal directions~\citep{spielman2011effective}.
These constructions illustrate several ways in which $\Delta_V$ can be specified.
Other fidelity criteria may be used when different relational properties are important.
